\begin{table}[!t]
\centering
\footnotesize
\caption{Practical guidelines for SFDA deployment in bearing fault diagnosis.}
\label{tab:guidelines}
\begin{tabular}{lllp{6cm}}
\toprule
ID & Guideline & Evidence & Rationale \\
\midrule
G1 & Report macro-F1 alongside accuracy & IV-A & Accuracy can overestimate performance by up to 33\,pp (NRC: 57.17 vs 27.39) \\
G2 & Disclose full implementation details & IV-E & Optimizer, freeze strategy, and training stages cause 7--37\,pp drift \\
G3 & Use Class Shift as default monitor & V-A & Pooled AUC = 0.853 across two datasets; best universal default \\
G4 & Calibrate thresholds with $\geq$20 normal runs & V-C & 10$\rightarrow$20 runs: Sens +0.20, Spec +0.30 \\
G5 & Select signal per method when possible & V-C, Tab.~\ref{tab:signal} & SHOT/RPSWD: Entropy; TENT: Class Shift \\
G6 & Do not rely on denoising alone & VI-A & Wavelet denoising: SNR +0.94\,dB, collapse unchanged (56.39\%) \\
G7 & Do not rely on monitoring-intervention alone & VI-A & Closed-loop intervention: $-$3.30\,pp, effective rate 35\% \\
G8 & Prefer conservative lr when no validation is possible & IV-B & SHOT lr=1e-4 retains 94.52\% at 0\,dB vs 58.80\% at lr=1e-3 \\
G9 & Audit method components before deployment & IV-F & RPSWD minus soft-weighting gains +11.31\,pp; components can be harmful \\
\bottomrule
\end{tabular}
\end{table}
